%---------------------------------------------------------------------
% UNIT 1 --- PROBABILITY THEORY
% Source: Section 1.9 Self Assessment Questions, pp. 25-26.
%---------------------------------------------------------------------
\chapter{Probability Theory}

\srcnote{Source: \emph{Business Mathematics} 5405/1429, \S1.9 \saq{Self Assessment Questions}, pp.~25--26.}

\begin{keyidea}[The four rules the whole unit runs on]
\begin{align*}
\text{Addition (general)}    &: && \prob{A\cup B}=\prob{A}+\prob{B}-\prob{A\cap B}\\
\text{Addition (exclusive)}  &: && \prob{A\cup B}=\prob{A}+\prob{B} \quad\text{when } A\cap B=\varnothing\\
\text{Multiplication (indep.)}&: && \prob{A\cap B}=\prob{A}\prob{B}\\
\text{Conditional}           &: && \prob{A\mid B}=\frac{\prob{A\cap B}}{\prob{B}},\qquad \prob{B}\neq 0
\end{align*}
Everything below is one of these four, or a complement $\prob{A'}=1-\prob{A}$.
\end{keyidea}

%---------------------------------------------------------------------
\section{Sample space of an odd draw}

\begin{question}{Q.1}
Find the sample space for choosing an odd number from 1 to 15 at random.
\end{question}

\begin{solution}
The sample space is the list of every outcome the experiment can produce.
Here the experiment is ``pick one odd number between 1 and 15'', so the
outcomes are the odd integers in that range:
\[
S=\{1,3,5,7,9,11,13,15\},\qquad n(S)=8 .
\]
Each is equally likely, so any single outcome has probability $1/8$.

\ans{$S=\{1,3,5,7,9,11,13,15\}$, with $n(S)=8$.}
\end{solution}

%---------------------------------------------------------------------
\section{Mutually exclusive versus collectively exhaustive}

\begin{question}{Q.2}
What is the difference between mutually exclusive events and collectively
exhaustive events?
\end{question}

\begin{solution}
These describe two different things: one is about \emph{overlap}, the other
about \emph{coverage}. They are independent properties, and a set of events can
have either, both, or neither.

\begin{center}
\small\sffamily
\begin{tabular}{@{}L{3.4cm} L{5.4cm} L{5.4cm}@{}}
\toprule
 & \textbf{Mutually exclusive} & \textbf{Collectively exhaustive} \\
\midrule
Condition & $A\cap B=\varnothing$ & $A\cup B=S$ \\
In probability & $\prob{A\cap B}=0$ & $\prob{A\cup B}=1$ \\
Plain reading & They cannot both happen. & At least one must happen. \\
Concerns & Overlap between events & Coverage of the sample space \\
\bottomrule
\end{tabular}
\end{center}

Take a single roll of a die, $S=\{1,2,3,4,5,6\}$:

\begin{itemize}
  \item $A=\{\text{even}\},\ B=\{\text{odd}\}$ --- both mutually exclusive
        \emph{and} collectively exhaustive. No overlap, and together they
        cover $S$.
  \item $A=\{1,2,3\},\ B=\{3,4,5,6\}$ --- collectively exhaustive but
        \emph{not} mutually exclusive: the outcome 3 is in both.
  \item $A=\{1\},\ B=\{2\}$ --- mutually exclusive but \emph{not}
        collectively exhaustive: a roll of 5 is in neither.
\end{itemize}

\ans{Mutually exclusive means the events cannot occur together
($A\cap B=\varnothing$). Collectively exhaustive means one of them must occur
($A\cup B=S$). The first is a statement about overlap, the second about
coverage; neither implies the other.}
\end{solution}

\begin{pitfall}
The simplified addition rule $\prob{A\cup B}=\prob{A}+\prob{B}$ is valid
\emph{only} for mutually exclusive events. Applying it to overlapping events
double-counts the intersection and is the single most common error in this unit.
\end{pitfall}

%---------------------------------------------------------------------
\section{Three independent admissions}

\begin{question}{Q.3}
The probability that an applicant for pilot school will be admitted is 0.5.
If three applicants are selected at random, what is the probability that
\textbf{(a)} all three will be admitted, \textbf{(b)} none will be admitted,
\textbf{(c)} only one will be admitted?
\end{question}

\begin{solution}
The three applicants are selected at random and independently, so
$\prob{\text{admit}}=0.5$ and $\prob{\text{reject}}=0.5$ for each, and
probabilities multiply across applicants.

\medskip
\textbf{(a) All three admitted.} One way this can happen, and the three events
are independent:
\[
\prob{3\text{ admitted}}=(0.5)(0.5)(0.5)=(0.5)^{3}=0.125 .
\]

\textbf{(b) None admitted.} Again one arrangement:
\[
\prob{0\text{ admitted}}=(0.5)^{3}=0.125 .
\]

\textbf{(c) Exactly one admitted.} Now the count matters. The admitted
applicant can be the first, the second or the third --- three arrangements, each
of probability $(0.5)(0.5)(0.5)=0.125$:
\[
\prob{1\text{ admitted}}=\binom{3}{1}(0.5)^{1}(0.5)^{2}=3(0.125)=0.375 .
\]

\vcheck{The four cases $0,1,2,3$ admitted must exhaust the sample space:
$0.125+0.375+0.375+0.125=1$. They do.}

\ans{(a) $0.125$ \quad (b) $0.125$ \quad (c) $0.375$}
\end{solution}

\begin{pitfall}
Part (c) is where marks go. ``Only one'' is not $(0.5)(0.5)^{2}$ --- you must
multiply by $\binom{3}{1}=3$ to count \emph{which} applicant it was. Parts (a)
and (b) need no such factor because there is only one way for all or none.
\end{pitfall}

%---------------------------------------------------------------------
\section{Grades: complements of exclusive events}

\begin{question}{Q.4}
A student estimates the probability of attaining an A in a maths course at 0.4
and the probability of attaining a B at 0.3. What is the probability that
he/she \textbf{(a)} will not receive an A, \textbf{(b)} will not receive a B,
\textbf{(c)} will receive neither an A nor a B?
\end{question}

\begin{solution}
Let $A$ be the event of an A grade and $B$ of a B grade, with $\prob{A}=0.4$
and $\prob{B}=0.3$. A single course yields a single grade, so $A$ and $B$ are
mutually exclusive: $\prob{A\cap B}=0$.

\medskip
\textbf{(a)} By the complement rule,
$\prob{A'}=1-\prob{A}=1-0.4=0.6$.

\textbf{(b)} $\prob{B'}=1-\prob{B}=1-0.3=0.7$.

\textbf{(c)} ``Neither A nor B'' is the complement of ``A or B''. Because the
two are mutually exclusive the addition rule loses its correction term:
\[
\prob{A\cup B}=\prob{A}+\prob{B}=0.4+0.3=0.7,
\]
\[
\prob{(A\cup B)'}=1-0.7=0.3 .
\]

\ans{(a) $0.6$ \quad (b) $0.7$ \quad (c) $0.3$}
\end{solution}

\begin{pitfall}
Part (c) is \emph{not} $\prob{A'}\times\prob{B'}=0.6\times0.7=0.42$. That
multiplication assumes $A'$ and $B'$ are independent, and they are not --- one
grade excludes the other. Use the complement of the union.
\end{pitfall}

%---------------------------------------------------------------------
\section{Urn of dotted and striped balls}

\begin{question}{Q.5}
An urn contains 8 green dotted balls, 10 green striped balls, 12 blue dotted
balls and 10 blue striped balls. If a ball is selected at random from the urn,
what is the probability that the ball will be \textbf{(a)} green or striped,
\textbf{(b)} dotted, \textbf{(c)} blue or dotted?
\end{question}

\begin{solution}
Lay the counts out as a two-way table first. Almost every mark in this question
is won or lost at this step.

\begin{center}
\small\sffamily
\begin{tabular}{@{}l C{2.2cm} C{2.2cm} C{2.2cm}@{}}
\toprule
 & \textbf{Dotted} & \textbf{Striped} & \textbf{Total} \\
\midrule
Green & 8  & 10 & \textbf{18} \\
Blue  & 12 & 10 & \textbf{22} \\
\midrule
\textbf{Total} & \textbf{20} & \textbf{20} & \textbf{40} \\
\bottomrule
\end{tabular}
\end{center}

\textbf{(a) Green or striped.} These overlap --- the 10 green striped balls are
in both --- so the general addition rule applies:
\[
\prob{G\cup S}=\prob{G}+\prob{S}-\prob{G\cap S}
=\frac{18}{40}+\frac{20}{40}-\frac{10}{40}=\frac{28}{40}=0.70 .
\]

\textbf{(b) Dotted.} Straight off the column total:
\[
\prob{D}=\frac{20}{40}=0.50 .
\]

\textbf{(c) Blue or dotted.} Overlap is the 12 blue dotted balls:
\[
\prob{B\cup D}=\frac{22}{40}+\frac{20}{40}-\frac{12}{40}=\frac{30}{40}=0.75 .
\]

\vcheck{Counting directly: green or striped $=8+10+10=28$ of 40, and blue or
dotted $=12+10+8=30$ of 40. Both agree.}

\ans{(a) $28/40=0.70$ \quad (b) $20/40=0.50$ \quad (c) $30/40=0.75$}
\end{solution}

%---------------------------------------------------------------------
\section{Four consecutive sixes}

\begin{question}{Q.6}
A single die is rolled and each side has an equal chance of occurring. What is
the probability of rolling four consecutive 6s?
\end{question}

\begin{solution}
Each roll is independent of the others --- a die has no memory --- and
$\prob{6}=\tfrac{1}{6}$ on every roll. Independence lets the probabilities
multiply:
\[
\prob{\text{four }6\text{s}}=\left(\frac{1}{6}\right)^{4}
=\frac{1}{1296}\approx 0.00077 .
\]

\ans{$\dfrac{1}{1296}\approx 0.00077$, i.e.\ about 8 times in every 10{,}000
attempts.}
\end{solution}

%---------------------------------------------------------------------
\section{Conditional probability in a mixed urn}

\begin{question}{Q.7}
A ball is selected at random from an urn containing three red striped balls,
eight solid \textup{[red]} balls, six yellow striped balls, four solid yellow
balls and four blue striped balls.
\textbf{(a)} What is the probability that the ball is yellow, given that it is
striped? \textbf{(b)} What is the probability that the ball is striped, given
that it is red? \textbf{(c)} What is the probability that the ball is blue,
given that it is solid coloured?
\end{question}

\srcnote{The book prints ``eight solid balls'' without a colour. The colour is
red: it is the only one of the three that would otherwise have no solid
entry, and parts (b) and (c) both depend on it. That reading is used here and
is flagged in the working.}

\begin{solution}
Tabulate, with a zero where the book lists nothing:

\begin{center}
\small\sffamily
\begin{tabular}{@{}l C{2.2cm} C{2.2cm} C{2.2cm}@{}}
\toprule
 & \textbf{Striped} & \textbf{Solid} & \textbf{Total} \\
\midrule
Red    & 3  & 8 & \textbf{11} \\
Yellow & 6  & 4 & \textbf{10} \\
Blue   & 4  & 0 & \textbf{4} \\
\midrule
\textbf{Total} & \textbf{13} & \textbf{12} & \textbf{25} \\
\bottomrule
\end{tabular}
\end{center}

A conditional probability restricts attention to one row or one column, and
renormalises within it.

\medskip
\textbf{(a) $\prob{\text{yellow}\mid\text{striped}}$.} The condition confines us
to the striped column, which holds 13 balls; 6 of them are yellow:
\[
\prob{Y\mid S}=\frac{\prob{Y\cap S}}{\prob{S}}
=\frac{6/25}{13/25}=\frac{6}{13}\approx 0.4615 .
\]

\textbf{(b) $\prob{\text{striped}\mid\text{red}}$.} Confine to the red row, 11
balls, of which 3 are striped:
\[
\prob{S\mid R}=\frac{3/25}{11/25}=\frac{3}{11}\approx 0.2727 .
\]

\textbf{(c) $\prob{\text{blue}\mid\text{solid}}$.} Confine to the solid column,
12 balls. None is blue:
\[
\prob{B\mid \text{solid}}=\frac{0/25}{12/25}=0 .
\]
This is a genuine answer, not a failure of the method: conditioning on ``solid''
makes ``blue'' impossible, because every blue ball in the urn is striped.

\ans{(a) $6/13\approx 0.462$ \quad (b) $3/11\approx 0.273$ \quad (c) $0$}
\end{solution}

\begin{pitfall}
The denominator of a conditional probability is the \emph{condition}, never the
grand total. Writing $\prob{Y\mid S}=6/25$ instead of $6/13$ throws away the
entire point of the question.
\end{pitfall}

%---------------------------------------------------------------------
\section{Recovering an intersection from a union}

\begin{question}{Q.8}
Suppose that $E$ and $F$ are events with $\prob{E}=0.2$, $\prob{F}=0.5$ and
$\prob{E\cup F}=0.6$. Determine
\textbf{(a)} $\prob{E\cap F}$, \textbf{(b)} $\prob{E\mid F}$,
\textbf{(c)} $\prob{F\mid E}$.
\end{question}

\begin{solution}
\textbf{(a)} Rearrange the general addition rule to make the intersection the
subject:
\[
\prob{E\cup F}=\prob{E}+\prob{F}-\prob{E\cap F}
\;\Longrightarrow\;
\prob{E\cap F}=\prob{E}+\prob{F}-\prob{E\cup F},
\]
\[
\prob{E\cap F}=0.2+0.5-0.6=0.1 .
\]

\textbf{(b)} \[
\prob{E\mid F}=\frac{\prob{E\cap F}}{\prob{F}}=\frac{0.1}{0.5}=0.2 .
\]

\textbf{(c)} \[
\prob{F\mid E}=\frac{\prob{E\cap F}}{\prob{E}}=\frac{0.1}{0.2}=0.5 .
\]

\medskip
\textbf{A remark worth a mark.} Notice that $\prob{E\mid F}=0.2=\prob{E}$ and
$\prob{F\mid E}=0.5=\prob{F}$: knowing one event tells you nothing about the
other. Equivalently $\prob{E}\prob{F}=(0.2)(0.5)=0.1=\prob{E\cap F}$. So $E$
and $F$ are \emph{independent}. If a follow-up part asks you to test
independence, this is the test.

\ans{(a) $0.1$ \quad (b) $0.2$ \quad (c) $0.5$; and $E,F$ are independent, since
$\prob{E\cap F}=\prob{E}\prob{F}$.}
\end{solution}

%---------------------------------------------------------------------
\section{Three kings without replacement}

\begin{question}{Q.9}
What is the probability of drawing three cards, without replacement, from a
deck of cards and getting three kings?
\end{question}

\begin{solution}
Without replacement means the deck shrinks and the count of kings shrinks with
it, so each factor must be recomputed:

\begin{center}
\small\sffamily
\begin{tabular}{@{}l C{2.6cm} C{2.6cm} C{2.4cm}@{}}
\toprule
\textbf{Draw} & \textbf{Kings left} & \textbf{Cards left} & \textbf{Probability} \\
\midrule
1st & 4 & 52 & $4/52$ \\
2nd & 3 & 51 & $3/51$ \\
3rd & 2 & 50 & $2/50$ \\
\bottomrule
\end{tabular}
\end{center}

By the multiplication rule for dependent events,
\[
\prob{3\text{ kings}}=\frac{4}{52}\cdot\frac{3}{51}\cdot\frac{2}{50}
=\frac{24}{132600}=\frac{1}{5525}\approx 0.000181 .
\]

\vcheck{By combinations, $\binom{4}{3}\big/\binom{52}{3}=4/22100=1/5525$. Same.}

\ans{$\dfrac{1}{5525}\approx 0.00018$}
\end{solution}

\begin{pitfall}
Using $\left(\tfrac{4}{52}\right)^{3}$ treats the draws as \emph{with}
replacement and gives $0.000455$ --- more than twice the true value. Read the
words ``without replacement'' as an instruction to change the denominator every
line.
\end{pitfall}

%---------------------------------------------------------------------
\section{Conditioning on stripes}

\begin{question}{Q.10}
An urn contains 18 red balls, 14 red striped balls, 16 yellow balls and 12
yellow striped balls.
\textbf{(a)} Given that a ball selected from the urn is striped, what is the
probability it is yellow?
\textbf{(b)} Given that a ball selected from the urn is not striped, what is the
probability it is red?
\end{question}

\begin{solution}
``18 red balls'' means plain red, as distinct from the 14 red striped:

\begin{center}
\small\sffamily
\begin{tabular}{@{}l C{2.4cm} C{2.4cm} C{2.2cm}@{}}
\toprule
 & \textbf{Striped} & \textbf{Not striped} & \textbf{Total} \\
\midrule
Red    & 14 & 18 & \textbf{32} \\
Yellow & 12 & 16 & \textbf{28} \\
\midrule
\textbf{Total} & \textbf{26} & \textbf{34} & \textbf{60} \\
\bottomrule
\end{tabular}
\end{center}

\textbf{(a)} Condition on the striped column (26 balls), of which 12 are yellow:
\[
\prob{Y\mid S}=\frac{12/60}{26/60}=\frac{12}{26}=\frac{6}{13}\approx 0.4615 .
\]

\textbf{(b)} Condition on the not-striped column (34 balls), of which 18 are red:
\[
\prob{R\mid S'}=\frac{18/60}{34/60}=\frac{18}{34}=\frac{9}{17}\approx 0.5294 .
\]

\ans{(a) $6/13\approx 0.462$ \quad (b) $9/17\approx 0.529$}
\end{solution}
